\documentclass[A4,11pt]{letter}
\pagestyle{empty}

\usepackage{amsmath}
\usepackage{amsfonts}
\usepackage{amssymb}
\usepackage{graphicx}
\usepackage{geometry}

\begin{document}

Dear Hossein Azizpour,

My name is Simon Sjöling, and I am applying for the doctoral student position in Generative Modeling, Super Resolution of Brain Imaging. I am currently finishing my master's thesis on real-time neural soft shadow rendering, which will be completed in June 2026, concluding five years of study at Chalmers University of Technology (M.Sc. in Data Science and AI, and B.Sc. in Electrical Engineering, Civilingenjör).

I am strongly interested in machine learning, especially when it comes to the engineering aspect of applying the technology to solve real problems. In my thesis, I explore how neural networks can approximate computationally expensive physical phenomena such as lighting and shadows for computer graphics, with an emphasis on accuracy and efficiency under strict real-time constraints ($> 60 \, \text{Hz}$).
A central part of this work has been iterating on the model and its implementation, using evaluations and ablation studies to identify which design choices have the highest impact on accuracy and performance. My main goal after my thesis is to contribute to the field of machine learning through research. Reading about this position made me excited to apply because improving the quality of brain MSI can have real impact on the understanding of brain processes, and the opportunity to develop new generative methods for super-resolution and imputation closely aligns with my goals.

My studies have given me a solid foundation in mathematics, physics, and computer science. During my master's, I have particularly focused on AI through coursework and project work which has included implementing, optimizing, and evaluating machine learning systems. I have worked with deep learning for vision, transformers, and generative methods such as diffusion models, and I am comfortable with modern ML workflows and libraries. Many of the courses I have taken, specifically image analysis and computer vision, have included building models specifically for medical applications, such as classifying medical imaging. I have found these projects to be particularly interesting, given the importance of the applications. Additionally, I have taken a course in biomedical engineering, and have a basic knowledge of medical imaging.

Much of my coursework has involved group projects, which I have enjoyed, and during my work at Ericsson I gained additional experience working in larger teams, taking responsibility for deliverables, and presenting results.

Beyond machine learning, I have also chosen courses mainly in mathematics, such as Non-linear Optimization, Partial Differential Equations, Stochastic Processes and Bayesian Inference, and financial mathematics, which have strengthened my mathematical understanding and advanced problem-solving skills. I have also taken High Performance Computing, which has been valuable for running computationally demanding tasks efficiently on both CPUs and, importantly for ML, GPUs.
\newpage
Outside of school, I have always had a passion for tinkering with software, which has resulted in numerous software implementations of RL agents, simple C++ and CUDA-based ML frameworks and more. This has strengthened my practical programming skills and my interest in the lower-level implementation of machine learning algorithms.

I am very excited at the prospect of contributing to the research of your group at KTH and learning from experts in the field. I believe that my background could be a good starting point, and I am highly motivated to develop through doctoral research with this group.

Best regards,\\
Simon Sjöling

\end{document}